\section{Threats to validity}
\label{sec:threats}

We argue each threat explicitly and quantify it where the data allow. Where a threat could not be
quantified, we say so rather than asserting that it is negligible.

\subsection{Survivorship and build attrition}

\textbf{Mechanism.} Repositories that build and run are not a random subset of the frame. Buildable
projects plausibly have simpler dependency graphs, fewer external services and smaller domains, and
these properties may correlate with both engineering care and oracle quality. The direction of the
resulting bias is not obvious: simpler domains are easier to test well (upward bias on score), while
projects with sophisticated suites are likelier to require services, browsers or databases and thus to
drop out (downward bias, because the strongest suites become unmeasurable).

\textbf{Quantification.} The attrition log records every candidate with a categorised reason, and the
funnel is reported in full (Section~\ref{sec:results}). The share of failures attributable to
\texttt{needs\_browser} and \texttt{missing\_service\_dependency} is the measurable footprint of the
downward mechanism. We report counts of \texttt{build\_fails}, \texttt{bdd\_fails} and
\texttt{no\_bdd\_execution} separately, precisely so that a reader can reconstruct which mechanism
dominates.

\textbf{Residual.} We cannot correct for an unmeasured common cause such as ``engineering care'' that
may drive both buildability and oracle quality. Bounds, not corrections, are available from
observational attrition data, and we do not claim otherwise.

\subsection{Mutant representativeness for business rules}

\textbf{Mechanism.} Generic code mutants are a \emph{proxy} for changes in business-rule behaviour, not
the same thing. A boundary mutant on a discount threshold plausibly models a rule change; a mutant in
a DTO getter does not. If the proxy is poor, the outcome measures something other than business-rule
sensitivity.

\textbf{Mitigation built into the design.} Three things limit the damage. Mutation is restricted to
the domain layer, so infrastructure and glue are excluded by construction. The operator set is
restricted to decision-logic operators. And ``business-rule relevance'' is treated as a separate,
explicitly labelled judgement rather than assumed.

\textbf{What we could not do.} The pre-registered plan called for two independent human raters to
label a random sample of mutants as business-rule-relevant or not, and to report Cohen's $\kappa$.
\textbf{No independent human raters were available for this study, so no $\kappa$ is reported.} We
regard this as a genuine limitation rather than a formality: it means the strength of the proxy
assumption is asserted on design grounds (domain restriction plus operator restriction) rather than
measured. A study that reports a $\kappa$ for this construct would materially strengthen the
evidence, and we recommend it to anyone replicating this work. We did not substitute an automated
labeller for a human rater, because an automated agreement statistic would measure the consistency of
a heuristic with itself and would misrepresent it as inter-rater reliability.

\subsection{Reliability of the judgement-coded indicators}

\textbf{Mechanism.} D2 (assertion weakness) and D7 (mirror steps) contain judgement. D2's weakening
patterns are implemented as explicit syntactic rules, but deciding whether a rule such as ``swallowed
failure'' applies is a coding decision; D7's full definition requires a rater to inspect the cited
domain method.

\textbf{What is reported.} D2 is implemented as an explicit decision procedure with published rules,
and its confidence is flagged per repository: where more than a quarter of step definitions cannot be
resolved, D2 is marked \texttt{low\_confidence} and excluded from the confirmatory test. D7 is
reported only as an \emph{automated census} of its loose conjunction, with no claim of validated
precision.

\textbf{Residual.} Without a second rater we cannot separate indicator noise from true variation. Any
association involving D2 or D7 should be read as provisional. This is stated in the results tables
themselves, not only here.

\subsection{Coverage confounding}

\textbf{Mechanism.} A suite that executes little of the domain has few opportunities to kill mutants,
so a low overall score can reflect reach rather than discrimination. The literature is explicit that
raw coverage--mutation association largely dissolves once restricted to covered code
\cite{Jain2023MindTheGap}.

\textbf{Mitigation.} Coverage conditioning is mandatory: $\msall$, $\mscov$ and coverage are reported
together in every table, and coverage is a covariate in every model. The primary outcome is the
covered-only score. We go further and state the interpretation rule verbatim in the method so that it
cannot be quietly abandoned when results are inconvenient.

\subsection{BDD attribution}

\textbf{Mechanism.} A suite may be strong for reasons unrelated to BDD: mature projects have good unit
tests, review culture and stable domains. Attributing a high mutation score to the Gherkin layer would
be a category error.

\textbf{Mitigation.} The differential $\dk$ is computed on a common mutant set and reported alongside
the absolute score; $\mathit{BDDshare}$ quantifies how much of the full suite's kill the specification
layer accounts for. The design is observational and \textbf{does not establish that BDD causes better
or worse oracles}. We state this in the introduction, here, and in the conclusion, because it is the
single easiest claim for a reader to over-infer.

\subsection{Two-snapshot measurement of drift}

\textbf{Mechanism.} Drift is a change over time. Measuring at HEAD and at one ancestor yields a
\emph{change between two snapshots}, which is not a trend and cannot distinguish erosion from a single
reorganisation or a bulk edit.

\textbf{Mitigation.} Where mutation is compared across commits, mutants are matched by file, operator
class and original-line content hash, and the change is reported only on the common set, and only when
that set covers at least half of the HEAD mutants; otherwise the repository is flagged
\texttt{low\_overlap}. We use the phrase ``change between two snapshots'' throughout and never
``trend'' or ``rate''.

\subsection{Public repositories versus industrial practice}

\textbf{Mechanism.} Public repositories differ from industrial practice in ways that matter: they have
no organisation-specific business rules, no auditors, often no sustained maintenance, and they are
self-selected for public visibility. Conclusions about organisational practice drawn from them are
correspondingly weak.

\textbf{Position.} This is why the paper frames its contribution as an instrument plus a population
finding rather than as a verdict on BDD. The population finding is itself bounded: it describes public
topic-tagged repositories, not all BDD practice.

\subsection{Toolchain heterogeneity}

\textbf{Mechanism.} PIT, Stryker and mutmut define, name and count mutants differently, so their scores
are not on a common axis. Kintis et al.\ found equivalent-mutant counts of 203, 94, 382 and 43 for the
same subjects across four Java tools, with no tool subsuming the others
\cite{Kintis2018HowEffectiveTools}. Our own harness produced a direct local instance: for the same
domain logic, PIT generated 19 mutants while Stryker generated 43.

\textbf{Mitigation.} Analysis is stratified by language and tool; no raw pooled correlation is
reported as a headline; where a pooled figure appears it is an average of within-stratum associations
with a cluster bootstrap, and is labelled as such. Operator sets are stated per tool.

\subsection{Construct terminology}

\textbf{Mechanism.} ``Specification drift'' is not an established term, and if a reader interprets it
as a received construct the paper will appear to claim more prior consensus than exists.

\textbf{Mitigation.} The term is introduced explicitly as a definition, its established relatives are
cited, and Section~\ref{sec:related} states that the terminological contribution is minor.

\subsection{Environment and reproducibility}

\textbf{Mechanism.} Measurements were taken on a single host without container isolation, in which the
Java toolchain is a single JDK version and mutation tools must be invoked outside the process sandbox.
A different environment could produce a different corpus because different repositories would build.

\textbf{Mitigation.} All tool versions are pinned and recorded, every measurement carries a UTC
timestamp, harness commit and seed in an append-only run log, and the screening and analysis scripts
are published. The vulnerability is not hidden: because the corpus is environment-dependent, we
report the full funnel so that a replication can compare its yield against ours rather than only its
final $n$.

